% A model through rsdag: front ends build the graph, transforms extend it,
% the tape lowers it, Adaptive serves the tape to the solver.
\documentclass[tikz,border=4pt]{standalone}
% The diagrams' shared preamble: Latin Modern, grey ink and one blue
% accent on a transparent page. make.py writes colors.tex from
% docs/style.py.
\usepackage[T1]{fontenc}
\usepackage{lmodern}
\usepackage{amsmath}
\usetikzlibrary{positioning,arrows.meta,fit,calc,backgrounds}
\input{colors}
\tikzset{
  every picture/.style={line width=0.7pt, color=ink, line join=round,
    >={Stealth[length=5pt,width=4pt]}, node distance=6mm and 10mm},
  % graph nodes: an operation, an input, an output's name
  op/.style={circle, draw, minimum size=8.5mm, inner sep=0.5pt},
  input/.style={op, fill=accent, draw=accent, text=white},
  output/.style={inner sep=1.5pt},
  % blocks: one line, or a title over a note
  box/.style={draw, minimum height=8mm, minimum width=21mm, inner sep=3pt,
    align=center},
  note/.style={font=\footnotesize},
  % a group frames its members; its title stands above it
  group/.style={draw, inner sep=3.5mm},
  title/.style={anchor=south west, inner sep=1.5pt, font=\small},
  % edges: data, state handed from the prolog, a derivative's own, a
  % branch not taken
  data/.style={->},
  state/.style={->, dashed, accent},
  derived/.style={dashed},
  faded/.style={opacity=0.35},
  lbl/.style={font=\footnotesize, inner sep=2pt},
  listing/.style={align=left, font=\ttfamily\small, inner sep=0pt},
}

\begin{document}
\begin{tikzpicture}[box/.append style={minimum width=23mm, minimum height=11mm}]
  \node[box, minimum width=17mm, minimum height=8mm] (py) at (0, 0) {Python};
  \node[box, minimum width=17mm, minimum height=8mm] (rs) at (0, 1.2) {Rust};
  \node[box, minimum width=17mm, minimum height=8mm] (mod) at (0, -1.2) {Module};
  \node[box, fill=accent, draw=accent, text=white] (graph) at (2.9, 0)
    {Graph\\[-2pt]\footnotesize hash-consed};
  \node[box, right=8mm of graph] (tr)
    {Transforms\\[-2pt]\footnotesize derivatives};
  \node[box, right=8mm of tr] (tape)
    {Tape\\[-2pt]\footnotesize prolog, main};
  \node[box, minimum height=8mm, minimum width=25mm, right=14mm of tape] (sp) {Specialized};
  \node[box, minimum height=8mm, minimum width=25mm, above=4mm of sp] (in) {Interpreter};
  \node[box, minimum height=8mm, minimum width=25mm, below=4mm of sp] (nat) {Native code};
  \node[group, fit=(in)(nat)] (ad) {};
  \node[title] at (ad.north west) {Adaptive};
  \node[box, right=8mm of ad, minimum width=17mm] (sol) {Solver};

  \coordinate (j) at ($(graph.west) + (-5mm, 0)$);
  \draw (rs.east) -| (j);
  \draw (mod.east) -| (j);
  \draw (py.east) -- (j);
  \draw[data] (j) -- (graph);
  \draw[data] (graph) -- (tr);
  \draw[data] (tr) -- (tape);
  \coordinate (s) at ($(tape.east) + (5mm, 0)$);
  \draw (tape.east) -- (s);
  \draw[data] (s) |- (in.west);
  \draw[data] (s) -- (sp.west);
  \draw[data] (s) |- (nat.west);
  \draw[data] (ad) -- (sol);
\end{tikzpicture}
\end{document}
